\section{Empirical strategy}\label{sec:strategy}

\subsection{Layered fixed-effects decomposition}

Let $R_i$ indicate a non-metro project. The paper estimates the sequence
\[
L_i \;=\; \alpha + \beta R_i + \delta_{t(i)} + \eta_{q(i)} + \mu_{s(i)}
+ \gamma_{c(i)} + \varepsilon_i,
\]
adding origination-year ($\delta$), QALICB-type ($\eta$), state ($\mu$),
and CDE ($\gamma$) fixed effects in turn. One feature of the sequence
deserves flagging at the outset, because it governs how the layered columns
should be read. The workhorse CDE specification replaces the state effects
of the preceding column rather than adding to them, since intermediary
identity already absorbs much of what state effects capture. Column 7 of
Table~\ref{tab:main} therefore reports the fully nested model carrying both
sets of effects, and I refer to it whenever a comparison across layers needs
to be strictly nested. Each layer removes a candidate
composition story. The CDE step is the workhorse: with $\gamma_{c(i)}$
included, $\bhat$ is identified from rural--urban comparisons \emph{within}
intermediary, so it measures the within-CDE rural penalty---the
market-structure component. The shrinkage of $\bhat$ upon adding CDE fixed
effects measures the between-CDE selection component. Standard errors in
CDE-fixed-effects specifications are clustered at the CDE level, the level
of the fixed effect and the economic decision-maker.

\subsection{Interpretation and limits}

The decomposition is descriptive in the sense of
selection-on-observables: within-CDE allocation of deals across rural and
urban tracts is not random, and no LATE interpretation of $\bhat$ is
claimed. What the design does establish is an accounting fact with economic
content: whether the aggregate rural penalty survives comparisons inside
the intermediary, where deal-structuring capacity, underwriting practice,
and investor relationships are held fixed. A causal companion
design---fuzzy RDD at the LIC eligibility threshold, run separately for
metro and non-metro samples---requires tract-level ACS covariates and is
deferred to the follow-on paper.

One further caution belongs here, because it constrains the language the
paper is entitled to use. The between-CDE component is a composition fact;
its ``selection'' reading is an interpretation. Rural market conditions
could in principle cause low-leverage intermediaries to specialize in rural
deployment, in which case part of the between-CDE component is market
structure operating through intermediary business models. The paper's
secure claim is the within-CDE null; the decomposition locates the gap, and
the discussion section is explicit about which readings of its location the
data can and cannot separate.

\subsection{Excess mass at the 20\% target}

For each CDE $j$, define its cumulative non-metro deployment share $s_j$,
computed both as a count share, $n_{j,\mathrm{nonmetro}}/n_{j,\mathrm{total}}$,
and as a dollar share of QLICI, since the requirement is written in terms of
investment. If the 20\% target binds at the deployment margin, optimizing
CDEs should pile up at $s_j = 0.20$ and produce excess mass in the sense of
\citet{chetty2011} and \citet{kleven2016}. The test compares empirical mass
in a $\pm 2.5$ percentage-point window around 20\% against a degree-3
polynomial counterfactual fitted outside the window, over the 310 CDEs with
at least five transactions, with a CDE-level bootstrap for inference. I
avoid describing the target as a notch. The public data reveal no
discontinuity in a CDE's objective or budget set at the threshold, and
whatever consequence attaches to falling short operates through future
allocation rounds that the release does not observe. Because the
proportionality instruction dates from 2006, I report the test for the full
period and again for origination years from 2007 onward.
